\section{Related Work}\label{sec: related work}

\AreaName touches on many areas of related work which we survey next.

\xhdr{Statistical Relational Learning} Since the foundation of the field of AI, sought to design systems capable of reasoning about entities and their relations, often by explicitly building graph structures \citep{minsky1974framework}. Each new era of AI research also brought its own form of relational learning. A prominent instance is statistical relational learning \citep{de2008logical}, a common form of which seeks to describe objects and relations in terms of first-order logic, fused with graphical models to model uncertainty \citep{getoor2001learning}. These descriptions can then be used to generate new ``knowledge'' through inductive logic programming \citep{lavrac1994inductive}. Markov logic networks, a prominent statistical relational approach, are defined by a collection of first-order logic formula with accompanying scalar weights \citep{richardson2006markov}. This information is then used to  define a probability distribution over possible worlds (via Markov random fields) which enables probabilistic reasoning about the truth of new formulae. We see \AreaName as inheriting this lineage, since both approaches operate on data with rich relational structure, and both approaches integrate relational structure into the model design. Of course, there are important distinctions between the two methods too, such as the natural scalability of graph neural network-based methods, and that \AreaName does not rely on first-order logic to describe data, allowing broad applicability to relations that are hard to fit into this form. 



\xhdr{Tabular Machine Learning} Tree based methods, notably XGBoost \citep{chen2016xgboost}, remain key workhorses of enterprise machine learning systems due to their scalability and reliability. In parallel, efforts to design deep learning architectures for tabular data have continued \citep{huang2020tabtransformer,arik2021tabnet,gorishniy2021revisiting,gorishniy2022embeddings,chen2023trompt}, but have struggled to clearly establish dominance over tree-based methods \citep{shwartz2022tabular}. The vast majority of tabular machine learning focuses on the single table setting, which we argue forgoes use of the rich interconnections between relational data. As such, it does not address the key problem, which is how to get the data from a multi-table to a single table representation. Recently, a nascent body of work has begun to consider multiple tables. For instance, \cite{zhu2023xtab} pre-train tabular Transformers that generalize to new tables with unseen columns. 

\xhdr{Knowledge Graph Embedding}  Knowledge graphs store relational data, and highly scalable knowledge graph embeddings methods have been developed over the last decade to embed data into spaces whose geometry reflects the relational struture~\citep{bordes2013translating,wang2014knowledge,wang2017knowledge}. Whilst also dealing with relational data, this literature differs from this present work in the task being solved. The key task of knowledge graph embeddings is to predict missing entities (Q: Who was Yann LeCun's postdoc advisor? A: Geoffrey Hinton)  or relations  (Q: Did Geoffrey Hinton win a Turing Award? A: Yes). To assist in such \emph{completion} tasks, knowledge graph methods learn an embedding space with the goal of exactly preserving the relation semantics between entities. This is different from \AreaName, which aims to make \emph{predictions} about entities, or groups of entities. Because of this, \AreaName seeks to leverage relations to learn entity representations, but does not need to learn an embedding space that perfectly preserves all relation semantics. This gives more freedom and flexibility to our models, which may discard certain relational information it finds unhelpful. Nonetheless, adopting ideas from knowledge graph embedding may yet be fruitful.
 
\xhdr{Deep Learning on Relational Data} Proposals to use message passing neural networks on relational data have occasionally surfaced within the research community. In particular, \cite{schlichtkrull2018relational, cvitkovic2020supervisedrd, vsir2021deep}, and \cite{zahradnik2023deep} make the connection between relational data and graph neural networks and explore it with different network architectures, such as heterogeneous message passing. However, our aim is to move beyond the conceptual level, and clearly establish deep learning on relational data as a subfield of machine learning. Accordingly, we focus on the components needed to establish this new area and attract broader interest: (1) a clearly scoped design space for neural network architectures on relational data, (2) a carefully chosen suite of benchmark databases and predictive tasks around which the community can center its efforts, (3) standardized data loading and splitting, so that temporal leakage does not contaminate experimental results, (4) recognizing time as a first-class citizen, integrated into all sections of the experimental pipeline, including temporal data splitting, time-based forecasting tasks, and temporal-based message passing, and (5) standardized evaluation protocols to ensure comparability between reported results.


